\begin{center}
{\fontsize{18}{22}\selectfont\bfseries\Titulo\par}
\vspace{10mm}
{\large\Autores\par}
\vspace{5mm}
{\small Artículo IV de la serie HMT--MD\par}
\vspace{2mm}
{\small Borrador de trabajo para revisión académica\par}
{\footnotesize Integración y recopilación documental generadas por ChatGPT Astra\par}
\end{center}
\vspace{6mm}
\begin{center}\textbf{Resumen}\end{center}
Se desarrolla la prolongación excepcional y dimensional de un núcleo formal
matemático denominado \emph{Holografía Modular Triádica}, en adelante HMT.
La construcción comienza en el soporte aritmético \emph{All Possible Paths}
(APP), incorpora la unidad ternaria orientada TRIT y compone sus operaciones
mediante el \emph{Topological Phase Kernel} (TPK). El estado resultante
conserva conjuntamente orientación, hojas aritméticas, residuos, cocientes,
frontera y memoria. Su prolongación compatible suministra la fuente común
de la realización numérica y de la incidencia excepcional.

El registro dodecafásico \(K\), construido dentro de esa cadena, cumple dos
funciones que se desarrollan mediante mapas explícitos. Su incidencia
interviene en la selección de una bandera de Witt y en la construcción
reticular que alcanza Leech y el álgebra de operadores de vértice
\(V^\natural\). Su proyección sobre una componente irreducible de la
representación de \(A_5\) determina una dirección no nula y la reducción
ortogonal \(12\to11\to10\). La construcción excepcional se acompaña de
los automorfismos de los sectores y de las realizaciones de orbifold de
órdenes \(2\) y \(3\), manteniendo identificados sus dominios.

La fase nonádica admite una representación de Weyl a cada profundidad.
Las inclusiones que conservan la fibra completa producen un álgebra
uniformemente hiperfinita y su realización tracial de tipo \(II_1\).
El intercambio de dos coordenadas enteras y la inversión del radio
determinan una dualidad exacta, tanto para la forma cuadrática compacta
como para su Hamiltoniano autoadjunto. Se prueba asimismo su formulación
covariante bajo el cambio de sección residual. Las pantallas reales,
combinatorias y reticulares se coordinan por sus morfismos, incluida la
reducción isotrópica \(26\to24\).

Finalmente, la composición de trayectorias con frontera conservada
define un cociclo de acción. La realización undecadimensional se formula
con sus campos, dominios operatorios y entrelazamientos: las identidades
algebraicas demostradas y el teorema de transporte de la supercarga
determinan con precisión el alcance de la interfaz física. El artículo
reúne en su propio cuerpo el núcleo común, la generación de los registros
y el corredor excepcional que necesitan estas construcciones.

\medskip
\noindent\textbf{Palabras clave:} Holografía Modular Triádica; estructura
discreta del continuo; registro dodecafásico; simetrías excepcionales;
Moonshine; dualidad T; representaciones de Weyl; pantallas dimensionales;
teoría M.
